\uuid{o4PH}
\exo7id{127}
\titre{exo7 127}
\auteur{cousquer}
\organisation{exo7}
\datecreate{2003-10-01}
\isIndication{false}
\isCorrection{false}
\chapitre{Logique, ensemble, raisonnement}
\sousChapitre{Ensemble}
\cpge{mpsi-mp2i.raisonnement-vocabulaire-ensembliste.b, pcsi.raisonnement-vocabulaire-ensembliste.b, ptsi.raisonnement-vocabulaire-ensembliste.b, mpsi-mp2i.raisonnement-vocabulaire-ensembliste.a, pcsi.raisonnement-vocabulaire-ensembliste.a, ptsi.raisonnement-vocabulaire-ensembliste.a, bcpst1.logique-ensembles-raisonnement.b, bcpst1.logique-ensembles-raisonnement.a}
\module{Algèbre}
\niveau{L1}
\difficulte{}

\contenu{
\question{
Montrer que si $F$ et $G$ sont des sous-ensembles de $E$~:
$$ (F \subset G \iff F \cup G = G )\quad\mbox{et} \quad(F \subset G \iff
\complement F \cup G = E).$$
En déduire que~: 
$$ (F \subset G \iff F \cap G = F)
\quad\mbox{et} \quad(F \subset G \iff F \cap\complement G = \emptyset).$$
}
}
